\documentclass[../../main2.tex]{subfiles}

\begin{document}

\chapter{预测与因果：能预测，就等价于有因果}
\label{ch:prediction}

\begin{motivation}
	上一章遗留的问题：人的目标只能反推，反推的结果只能是概率断言。但概率断言到处都有——星座运势也给断言。凭什么说有的预测\textbf{站得住}？导论钉下的枢纽论断「能预测，就等价于有因果」，到现在还只是一句话。\\
	\textbf{核心动机}：把「预测」这个词收窄到能负责的范围——预测「改变世界之后，世界怎么回应」——并证明这种预测离开因果寸步难行。
\end{motivation}

\begin{logicmap}
\begin{tikzpicture}[node distance=0.85cm, every node/.style={draw, rectangle, rounded corners, align=center, font=\small}]
\node (q) {预测 = 什么？};
\node (c) [below=of q] {两种预测：趋势 vs 干预};
\node (d) [below=of c] {枢纽：能预测，等价于有因果};
\node (p) [below=of d] {方向：从因推果};
\node (n) [below=of p] {相关为什么不等于原因（抛给第 7 章）};
\draw[-{Stealth}] (q) -- (c);
\draw[-{Stealth}] (c) -- (d);
\draw[-{Stealth}] (d) -- (p);
\draw[-{Stealth}] (p) -- (n);
\end{tikzpicture}
\end{logicmap}

\section{两种预测，不是一个物种}

\begin{readingnote}
	你有没有想过，为什么保险公司从不预言谁哪天死，却是全人类最会算「死」的公司？\\
	因为它们做的不是预言，是\textbf{概率断言}——而且是带因果的那种。
\end{readingnote}

先把「预测」拆成两种，拆完你会立刻认出谁靠谱：

\begin{itemize}
	\item \textbf{趋势预测}：给定现在，推下一步。天气预报、潮汐表、地铁到站时间。世界不变，外推就灵。
	\item \textbf{干预预测}（interventional prediction）：问「如果我动一下这个世界，会怎样」。人工降雨会不会下雨、限购会不会压住房价、换专业会不会更好。
\end{itemize}

前一种只需要拟合历史，后一种必须知道\textbf{谁在起作用}。本书从头到尾说的「预测」，都是后一种。

\begin{tcolorbox}[colback=green!5, title=\textit{生活类比}]
	导航软件是趋势预测的顶峰：照着地图和车流外推。但你问它「如果我今天不开车，路上会少几辆车」——它闭嘴了。前者是算路，后者是算\textbf{你对世界的因果}。两件事中间隔着的，就是整本书。
\end{tcolorbox}

当下最让人困惑的几个问题，其实全是干预预测：

\begin{itemize}
	\item \textbf{房价}：限购、降息、学区改革，哪个真能压住价格？——「如果不限购，房价会更高吗？」
	\item \textbf{短视频}：未成年人模式上线，孩子就会早睡吗？——「如果没有这个功能，睡眠会差多少？」
	\item \textbf{高考}：新题型改革，刷旧题的人吃亏吗？——「如果没改革，名次会怎么变？」
	\item \textbf{流行病}：勤洗手真的少生病吗？——「如果没有这条倡议，感染率会高几成？」
\end{itemize}

四句话长得不一样，动作是同一个：先在脑子里造一个「没有 X 的平行宇宙」，再对比两个宇宙。这是预测的最小单元，叫\textbf{反事实推演}。

\begin{questionbox}
	\textit{现在你可能想反驳：「趋势预测不是也天天用吗？看 K 线炒股、看销量进货，不都在外推？」}\\
	\textbf{对，而且平稳年代外推挺好用。}但外推回答不了「该不该动手」。看 K 线能告诉你「照这个走势明天大概率涨」，回答不了「如果我今天清仓，会怎样」。而你人生中真正要做的决定，全是干预——换工作、开口表白、给孩子转学。做决定就要干预预测。
\end{questionbox}

\paragraph*{逻辑衔接}
	两种预测分家了：趋势靠外推，干预靠因果。但凭什么说后者「等价于」有因果？这句话值得当面证明——下一节拆枢纽。

\section{枢纽：能预测，就等价于有因果}

\begin{readingnote}
	你有没有想过，「吸烟导致肺癌」这个结论，是先有数据还是先有干预？\\
	1950 年代数据早就够多了，争论却拖了几十年——缺的那块拼图，正是\textbf{干预}。
\end{readingnote}

拆这句枢纽，两头各走一步。

\textbf{第一步：能预测「改变 A 之后的结果」，说明你知道 A 是不是原因。}如果你能预测「放开限购后房价会涨」，你其实已经断言了「限购对房价有作用」。预测得越准，这个断言越硬。

\textbf{第二步：知道 A 是原因，你就能预测「改变 A 之后的结果」。}因果知识的全部价值就在这——它保证「动手之后」的世界不是乱数。不知道因果，动手就是赌博。

\begin{tcolorbox}[colback=green!5, title=\textit{生活类比}]
	老中医和老菜农的区别。老菜农说「霜降之后的白菜更甜」——他知道\textbf{霜}是原因，所以他能预测：今年霜来得早，甜得早。如果他只记得「每年这时候白菜甜」，霜不来，他就懵了。知其然（规律）和知其所以然（因果），在变天那天高下立判。
\end{tcolorbox}

这就是寻找真正的幕后推手（causal inference，因果推断）的意思：不问「什么和什么一起出现」，问「谁在真正起作用」。三个历史现场，都在争这一件事：

\begin{itemize}
	\item \textbf{烟草与肺癌}：相关性 1950 年代就铁证如山，反对方仍说「体质同时导致吸烟和肺癌」。最后靠戒烟干预、剂量反应、跨国对照，因果才定案。
	\item \textbf{产褥热}：塞麦尔维斯发现「医生接生比助产士接生死亡率高」，推出原因是\textbf{医生的手}。他要求洗手，死亡率暴跌——干预验证了因果。他没能说清机制，被同行嘲讽到死。
	\item \textbf{坏血病}：水手吃柠檬不得坏血病——林德做了对照实验（干预！），因果才从「相关」升级为「可控」。英国海军靠这条多打赢了几场海战。
\end{itemize}

\lowfreq{诚实标注：现实中「完全预测得准」不存在，本书说的等价是在「同一理想化模型内」的等价。真实社会有噪声——这正是后面失败清单要处理的。}

\begin{questionbox}
	\textit{现在你可能想反驳：「天气预报不靠因果，不也预测得很准？」}\\
	\textbf{第二遍问这个问题的人，都是聪明人。}天气预报是趋势预测的王者，但它是「算路」，不是「算你」。它答不了「如果我人工降雨」——那要因果。再补一刀：天气预报能做到的极限，恰恰是「不改变天气」这个前提；一旦你想动手，它就退到因果问题后面去了。
\end{questionbox}

\begin{reader_anticipation}
	\item 相关不等于原因，这句话谁都会背——怎么\textbf{操作}？→ 第 7 章给你三把镊子：找混杂、造自然实验、看剂量反应。
	\item 因果知道了，「各占多少」怎么算？→ 第 8 章进法庭，看人类被迫算了上千年的那套手艺。
\end{reader_anticipation}

\paragraph*{逻辑衔接}
	枢纽立住了：干预预测离开因果寸步难行。但因果知识自己从哪来？下一节先解决一个方向问题——因果的箭头，该朝哪边读。

\section{方向：从因推果}

\begin{readingnote}
	你有没有想过，医生看病是「由果推因」——从症状倒推病因；政府评估政策却是「由因推果」——从政策推效果。两种都叫因果，凭什么书里只认一种？
\end{readingnote}

两个方向，各有地盘：

\begin{itemize}
	\item \textbf{果推因}（诊断）：病人发烧了，是什么病？医生列候选、查证据。这是「从结果猜原因」。
	\item \textbf{因推果}（预测）：如果开这刀、上这个药、上这条政策，会怎样？这是「从原因推结果」。
\end{itemize}

诊断给出的是\textbf{候选清单}——第 1 章说过，同一个结果可以来自不同的原因组合，倒推没有唯一解。要拍板，必须回到正向：对比「有它没它」两种世界。所以本书的因果箭头只朝一个方向：\textbf{从因推果}。倒推可以当线索，不能当判决。

\begin{tcolorbox}[colback=green!5, title=\textit{生活类比}]
	破案和审判是两回事。侦探由果推因，列嫌疑人（候选）；检察官必须由因推果，证明「没有他，案子就不会这样发生」。福尔摩斯可以猜，法官不许猜。本书学的是法官。
\end{tcolorbox}

历史学的限度也在这。历史学几乎全是果推因：从结果倒推「为什么罗马亡了」。它成立，但不完备——它产出候选，产不出「各占多少」。要分账，就得像下一章那样，造平行宇宙来对比。

\paragraph*{逻辑衔接}
	方向定了：从因推果。可现实里，原因从不排队——一堆东西一起变，谁是推手谁是路过的？「相关」这个最常见的陷阱，下一章正面拆。

\begin{thinkbox}
\begin{enumerate}
	\item 挑一个你正在纠结的决定（转专业、跳槽、搬城市）：写出你要预测的「干预」——「如果我做了 X，会怎样」。再写出你现在的依据：哪些是外推（趋势），哪些真摸到了因果？
	\item \textbf{反事实}：如果 1854 年伦敦霍乱爆发时，斯诺医生没有去标病例地图、没有拆掉那个水泵把手，「水传播霍乱」的因果要晚多少年才定案？没有这次干预，后世的防疫会少哪几条？
	\item \textbf{反事实}：如果当年烟草公司说的「体质论」是对的——吸烟和肺癌真是同一个体质的两个表现——世界会有什么不一样？戒烟倡议还会降低肺癌吗？用「干预预测」的标准检验一下这个假说。
\end{enumerate}
\end{thinkbox}

\end{document}
